%% 5-18 まとめ：微分したときのふるまい（第11回：ページ制限なし。表の1行ずつに「なぜ」と図を書き足した）
\newcommand\bkvromidashi[1]{\par\Needspace{6\baselineskip}\noindent\textbf{#1}\par}
\bkchapter{5}{18}{まとめ：微分したときのふるまい}{\bklevel{3}}{まとめ：微分したときのふるまい}
\bkmichishirube{5-14,5-15,5-17}{6-1,6-4,6-6}

\begin{goal}
\begin{itemize}[label={\color{bkTheme}●}, leftmargin=1.6\zw, itemsep=0pt, topsep=1pt]
\item \bklevel{3}\enspace 定数・$t^n$・$e^{at}$・$\sin\omega t$・$\cos\omega t$・$\ln t$ の微分と積分を、表にまとめて使える。
\item \bklevel{3}\enspace 「微分すると自分の $a$ 倍」「2回微分すると自分の $-\omega^2$ 倍」になる関数を答えられる。
\item \bklevel{3}\enspace はじめの状態から、式の中の数を決められる。
\end{itemize}
\end{goal}

\begin{hatena}
\noindent\begin{minipage}[c]{0.58\linewidth}
第6巻では、「変わる速さが、いまの状態で決まる」という式から、これからの動きを求めます。たとえば、
\begin{itemize}[leftmargin=1.2\zw, itemsep=0pt, topsep=2pt, label={\color{bkYellow!70!black}●}]
\item 微分すると、自分の $-2$ 倍になる関数は？
\item 2回微分すると、自分の $-9$ 倍になる関数は？
\end{itemize}
こんななぞなぞに答えるには、どんな関数を知っていればよいでしょう。
\end{minipage}\hfill
\begin{minipage}[c]{0.38\linewidth}
\centering
\begin{tikzpicture}[x=1mm, y=1mm]
  \foreach \i/\c in {0/bkBlue,1/bkOrange,2/bkPurple,3/bkTheme} {\draw[\c, line width=1pt, rounded corners=2pt, fill=\c!10] (\i*11,0) rectangle ++(9,13); \node[font=\sffamily\bfseries\small, text=\c] at (\i*11+4.5,6.5) {？};}
  \node[font=\sffamily\scriptsize, anchor=north] at (20,-2) {答えの候補のカード};
\end{tikzpicture}
\end{minipage}
\end{hatena}

\clearpage
第5巻で、微分と積分の道具と、$e^{at}$・$\sin$・$\cos$ という大切な関数がそろいました。この章では、それらを「微分したときのふるまい」で整理し、第6巻で使う「答えの候補リスト」を作ります。

%% ---------------- やってみよう・思い出そう ----------------
\begin{yattemiyou}[関数あてクイズ]
左のせつめいに合う関数を、右から選びましょう（ノートに線で結ぶ）。答えは微分して確かめます。\par\vspace{7pt}
\begin{bkcenter}
\begin{tblr}{colspec={Q[l,m]Q[c,m]Q[l,m]}, hline{1,Z}={0.6pt,bkGray}, rowsep=3pt, colsep=6pt}
（1）微分すると0 & & ア\enspace $e^{3t}$ \\
（2）微分すると、いつも2 & & イ\enspace $\sin2t$ \\
（3）微分すると、自分の3倍 & & ウ\enspace 5（定数） \\
（4）2回微分すると、自分の $-4$ 倍 & & エ\enspace $\ln t$ \\
（5）微分すると $\dfrac{1}{t}$ & & オ\enspace $2t+1$ \\
\end{tblr}
\end{bkcenter}
\end{yattemiyou}

\noindent 答えは（1）ウ、（2）オ、（3）ア、（4）イ、（5）エです。1つずつ微分して確かめます。
\begin{bkcenter}
\begin{bktable}{colspec={Q[l,m,wd=14mm] Q[l,m,wd=62mm] Q[l,m,wd=36mm]}, rowsep=1.5pt}
 & 微分して確かめる & 使ったきまり \\
（1）ウ & $(5)'=0$ & 定数の微分は0（5-4） \\
（2）オ & $(2t+1)'=2+0=2$ & $(t)'=1$（5-4） \\
（3）ア & $(e^{3t})'=3e^{3t}$。自分（$e^{3t}$）の3倍 & $(e^{at})'=ae^{at}$（5-14） \\
（4）イ & $(\sin2t)'=2\cos2t$、$(2\cos2t)'=-4\sin2t$ & 5-15 \\
（5）エ & $(\ln t)'=\dfrac{1}{t}$ & 5-14 \\
\end{bktable}
\end{bkcenter}
（4）は $\cos2t$ でも正解です（$(\cos2t)''=-4\cos2t$）。このように「せつめい（微分したときのふるまい）」から関数を当てることが、第6巻の微分方程式を解くことそのものです。

\begin{omoidasou}
\bkrecall{5-4}{$(t^n)'=nt^{n-1}$。定数の微分は0。定数倍は前に出せる。}
\bkrecall{5-13}{積分すると $+C$ がつく。$C$ は、はじめの状態（初期条件）で決まる。}
\bkrecall{5-14}{$(e^{at})'=ae^{at}$、$e^0=1$、$(\ln t)'=\dfrac{1}{t}$。$e^x=y$ なら $x=\ln y$}
\bkrecall{5-15}{$(\sin\omega t)'=\omega\cos\omega t$、$(\cos\omega t)'=-\omega\sin\omega t$。2回微分で $-\omega^2$ 倍}
\bkrecall{5-17}{テイラー展開。$e^t=1+t+\frac{t^2}{2!}+\frac{t^3}{3!}+\cdots$}
\end{omoidasou}

新しい計算はほとんどありません。これまでの章をふり返りながら、1つずつ確かめていきましょう。

%% ---------------- 数学の道具 ----------------
\bksection{doguu}{数学の道具}

\bksubsection{1}{微分と積分のふるまい表}

\label{bk:5-18:hyo}第5巻で出てきた関数を、1つの表にまとめます。$C$ は積分定数です（5-11）。
\begin{bkcenter}
\small
\begin{tblr}{colspec={Q[c,m]Q[c,m]Q[c,m]Q[c,m]Q[l,m]}, row{1}={bg=bkTheme!12, font=\sffamily\small}, hline{1,2,Z}={0.6pt,bkGray}, hline{3-7}={0.3pt,bkGray}, rowsep=2pt, colsep=5pt}
関数 $f(t)$ & 微分 $f'(t)$ & 2回微分 $f''(t)$ & 積分 $\int f\,\dd t$ & ふるまい \\
$c$（定数） & 0 & 0 & $ct+C$ & 変わらない \\
$t^n$ & $nt^{n-1}$ & $n(n-1)t^{n-2}$ & $\dfrac{t^{n+1}}{n+1}+C$ & 次数が1つ下がる \\
$e^{at}$ & $ae^{at}$ & $a^2e^{at}$ & $\dfrac{1}{a}e^{at}+C$ & 自分の $a$ 倍 \\
$\sin\omega t$ & $\omega\cos\omega t$ & $-\omega^2\sin\omega t$ & $-\dfrac{1}{\omega}\cos\omega t+C$ & $\frac{1}{4}$ 周期進む \\
$\cos\omega t$ & $-\omega\sin\omega t$ & $-\omega^2\cos\omega t$ & $\dfrac{1}{\omega}\sin\omega t+C$ & $\frac{1}{4}$ 周期進む \\
$\ln t$（$t>0$） & $\dfrac{1}{t}$ & $-\dfrac{1}{t^2}$ & $t\ln t-t+C$ & ゆっくりふえる \\
\end{tblr}
\end{bkcenter}
表の1行ずつ、どこで作ったかと、なぜそうなるかを短くふり返ります。
\begin{itemize}[leftmargin=1.2\zw, itemsep=2pt, topsep=2pt, label={\color{bkTheme}●}]
\item \textbf{定数}\quad 変わらない量は、変わる速さが0（5-4）。積分は、高さ $c$ の長方形の面積 $ct$（5-9）。
\item \textbf{$t^n$}\quad 微分は5-4の表から見つけた「指数を前に出して1へらす」。2回微分は、それをもう1回使ったもの（$nt^{n-1}$ の指数 $n-1$ を前に出す）。積分は逆に「指数を1ふやして、その数でわる」（5-11）。
\item \textbf{$e^{at}$}\quad 微分するたびに $a$ が1つかかる（5-14）ので、2回微分では $a\times a=a^2$。積分では逆に $a$ でわる。
\item \textbf{$\sin\omega t$・$\cos\omega t$}\quad 円運動の速度から作った（5-15）。微分するたびに $\omega$ がかかり、$\frac{1}{4}$ 周期ずつ進む。2回微分で $\omega^2$ がかかり、$-$ がつく。
\item \textbf{$\ln t$}\quad $e^t$ のグラフを折り返して微分した（5-14）。積分は部分積分で求めた（5-16）。2回微分は、$\dfrac{1}{t}=t^{-1}$ を $t^n$ の行で $n=-1$ として、$-1\times t^{-2}=-\dfrac{1}{t^2}$。
\end{itemize}
積や合成関数は、積の微分・合成関数の微分（5-6）、置換積分・部分積分（5-16）で、この表の形にもちこみます。

\begin{bkfigure}
\begin{tikzpicture}
\begin{scope}[x=1mm, y=1mm]
  \foreach \i/\ttl/\ax in {0/{$e^{0.4t}$：ふえる}/0, 1/{$e^{-t}$：へる}/0, 2/{$\sin t$：ゆれる}/10, 3/{$\ln t$：ゆっくりふえる}/6} {
    \begin{scope}[xshift=\i*33mm]
      \draw[bkGray, line width=0.6pt, -{Stealth[length=3pt]}] (0,0) -- (0,20);
      \draw[bkGray, line width=0.6pt, -{Stealth[length=3pt]}] (0,\ax) -- (27,\ax) node[anchor=west, font=\sffamily\scriptsize] {$t$};
      \node[font=\sffamily\footnotesize, anchor=north] at (13,-2) {\ttl};
    \end{scope}}
  \draw[bkBlue, line width=1.4pt] plot[domain=0:3.6, samples=30] ({7*\x}, {4*exp(0.4*\x)});
  \draw[bkBlue, line width=1.4pt, xshift=33mm] plot[domain=0:3.6, samples=30] ({7*\x}, {16*exp(-\x)});
  \draw[bkBlue, line width=1.4pt, xshift=66mm] plot[domain=0:6.283, samples=50] ({3.8*\x}, {10+7*sin(\x r)});
  \draw[bkBlue, line width=1.4pt, xshift=99mm] plot[domain=0.35:10, samples=50] ({2.4*\x}, {6+4.6*ln(\x)});
\end{scope}
\end{tikzpicture}
\bkfigcaption{4つの「ふるまい」。ふえる・へる・ゆれる・ゆっくりふえる。自然や機械の変化の多くは、この形か、その組み合わせになる。}
\end{bkfigure}

\bksubsection{2}{なぞなぞの答え：自分にもどる関数}

\label{bk:5-18:modoru}表の中で、微分しても「もとの形」にもどってくる関数が、とくに大切です。
\begin{bkcenter}
\small
\begin{tblr}{width=\linewidth, colspec={X[l,m]Q[c,m]Q[c,m]Q[l,m]}, row{1}={bg=bkTheme!12, font=\sffamily\small}, hline{1,2,Z}={0.6pt,bkGray}, hline{3-5}={0.3pt,bkGray}, rowsep=3pt, colsep=5pt}
なぞなぞ & 式 & 答えの関数 & 例 \\
微分すると自分の $a$ 倍 & $f'=af$ & $Ce^{at}$ & ふえる・へる（6-4） \\
2回微分で自分の $-\omega^2$ 倍 & $f''=-\omega^2f$ & $A\sin\omega t+B\cos\omega t$ & ゆれる（6-6） \\
2回微分で自分の $a^2$ 倍 & $f''=a^2f$ & $Ae^{at}+Be^{-at}$ & たおれる \\
微分すると定数 $v$ & $f'=v$ & $vt+C$ & 等速（6-3） \\
\end{tblr}
\end{bkcenter}
$C$、$A$、$B$ はどんな数でもかまいません。定数倍は、微分しても前にそのまま残るからです（5-4）。表の答えが正しいことを、1行ずつ微分して確かめます。
\begin{itemize}[leftmargin=1.2\zw, itemsep=2pt, topsep=2pt, label={\color{bkTheme}●}]
\item $f=Ce^{at}$ なら $f'=C\times ae^{at}=a\times(Ce^{at})=af$。
\item $f=Ae^{at}+Be^{-at}$ なら、
\[ f'=aAe^{at}-aBe^{-at},\qquad f''=a^2Ae^{at}+a^2Be^{-at}=a^2f \]
$e^{-at}$ は1回微分するごとに $-a$ がかかり、2回で $(-a)\times(-a)=a^2$ になるからです。
\item $f=vt+C$ なら $f'=v$。
\end{itemize}
（$A\sin\omega t+B\cos\omega t$ は、このあとの「2つをたしても答えになる」で確かめます。）「たおれる」は、手のひらに立てたほうきのように、少しかたむくと、かたむきがどんどん大きくなる動きです。$e^{at}$ の部分が、時間とともにどんどん大きくなります。

\bkvromidashi{\bkhatenaref の答え}
\begin{itemize}[leftmargin=1.2\zw, itemsep=2pt, topsep=2pt, label={\color{bkTheme}●}]
\item 「微分すると、自分の $-2$ 倍」は $f'=-2f$。表の1行めで $a=-2$ として、$f=Ce^{-2t}$。
\item 「2回微分すると、自分の $-9$ 倍」は $f''=-9f$。表の2行めとくらべると $\omega^2=9$。2乗して9になる正の数は3なので $\omega=3$。答えは $f=A\sin3t+B\cos3t$。
\end{itemize}
$\omega=9$ ではなく $\omega=3$ になることに注意しましょう。$\sin3t$ を2回微分すると、3が2回かかって $3\times3=9$ になるからです。

\bkvromidashi{2つをたしても答えになる}
$\sin3t$ と $\cos3t$ は、どちらも $f''=-9f$ をみたします。では、たしたもの $f=A\sin3t+B\cos3t$ はどうでしょう。2回微分すると、
\[ f''=A(-9\sin3t)+B(-9\cos3t)=-9(A\sin3t+B\cos3t)=-9f \]
で、たしたものも答えになります。$-9$ を前にくくり出せたのは、どちらの項にも同じ $-9$ がかかったからです。このように「答えどうしをたしても答え」になり、「答えを何倍しても答え」になる式を\textbf{線形}な式といいます。第6巻・第9巻の中心になる考えです。\bkidx{せんけい}{線形}

\noindent 線形でない式では、こうはなりません。たとえば $f'=f^2$（変わる速さが自分の2乗）では、$f$ が答えでも $2f$ は答えになりません（$2f$ を微分すると $2f'=2f^2$ ですが、$(2f)^2=4f^2$ なので合わない）。

\bksubsection{3}{数を決めるのは、はじめの状態}

\label{bk:5-18:shoki}$Ce^{at}$ の $C$ や、$A\sin\omega t+B\cos\omega t$ の $A$、$B$ は、なぞなぞだけでは決まりません。5-13と同じように、はじめの状態（初期条件）で決めます。

\begin{bkfigure}
\begin{tikzpicture}
\begin{axis}[bkplot, width=80mm, height=50mm, xmin=0, xmax=12.8, ymin=0, ymax=135, xtick={0,5,10}, ytick={40,80,120}, xlabel={$t$}, ylabel={$f$}]
  \addplot[bkthin, domain=0:12.5, samples=40] {120*exp(-0.1*x)};
  \addplot[bkcurve, domain=0:12.5, samples=40] {80*exp(-0.1*x)};
  \addplot[bkthin, domain=0:12.5, samples=40] {40*exp(-0.1*x)};
  \addplot[only marks, mark=*, mark size=1.8pt, bkOrange] coordinates {(0,80)};
  \draw[bkGray, line width=0.4pt] (axis cs:12.6,34.4) -- (axis cs:13.5,54);
  \draw[bkGray, line width=0.4pt] (axis cs:12.6,22.9) -- (axis cs:13.5,36);
  \draw[bkGray, line width=0.4pt] (axis cs:12.6,11.5) -- (axis cs:13.5,18);
  \node[font=\sffamily\scriptsize, anchor=west, text=bkGray] at (axis cs:13.6,54) {$C=120$};
  \node[font=\sffamily\scriptsize, anchor=west, text=bkBlue] at (axis cs:13.6,36) {$C=80$};
  \node[font=\sffamily\scriptsize, anchor=west, text=bkGray] at (axis cs:13.6,18) {$C=40$};
\end{axis}
\end{tikzpicture}
\bkfigcaption{3本の曲線は、どれも $f'=-0.1f$ をみたす $f=Ce^{-0.1t}$。$C$ がちがうと、スタートの高さがちがう。「$t=0$ で80」と決めると、オレンジの点を通る1本（青）だけが選ばれる。}
\end{bkfigure}

\bkvromidashi{例1}
$f'=-0.1f$、$f(0)=80$ のとき。答えの形は $f=Ce^{-0.1t}$ です。$t=0$ を入れると、$e^0=1$ なので、
\[ f(0)=Ce^{-0.1\times0}=Ce^0=C\times1=C\qquad\text{これが80なので}\quad C=80 \]
よって $f=80e^{-0.1t}$。$C$ は、ちょうど「はじめの値」になります。

\bkvromidashi{例2}
$\ddot{x}=-4x$、$x(0)=0.1$、$\dot{x}(0)=0$ のとき（ばねを0.1だけのばして、静かに手をはなした）。$-\omega^2=-4$ から $\omega=2$ なので、答えの形は $x=A\sin2t+B\cos2t$ です。
\begin{enumerate}[label={\bkcircled[bkTheme]{\arabic*}}, leftmargin=2.4\zw, itemsep=1pt, topsep=2pt]
\item はじめの位置：$t=0$ を入れる。$\sin0=0$、$\cos0=1$ なので、
\[ x(0)=A\times0+B\times1=B \]
これが0.1なので $B=0.1$。
\item はじめの速度：まず微分して $\dot{x}=2A\cos2t-2B\sin2t$。$t=0$ を入れると、
\[ \dot{x}(0)=2A\times1-2B\times0=2A \]
これが0なので $A=0$。
\end{enumerate}
だから $x=0.1\cos2t$ です。はしから手をはなすと、$\cos$ の形（山から始まる波）でゆれます。

\noindent 1回微分の式では数が1つ（$C$）、2回微分の式では数が2つ（$A$、$B$）。1回積分するごとに $C$ が1つ出てくる（5-13）ので、2回微分の式には、はじめの「位置」と「速度」の2つが必要です。この2つで、ばねの動きはただ1つに決まります。

\bksubsection{4}{ずらしても使える}

\label{bk:5-18:zurasu}お湯が冷めるとき、冷める速さは「部屋の温度との差」に比例します（6-4）。部屋が20\,℃なら、比例定数を $k$ として $T'=-k(T-20)$。このままでは表にない形です。

そこで、部屋の温度との差 $u=T-20$ を主役にします。$T=80$ なら $u=60$、$T=50$ なら $u=30$ です。
\begin{bkfigure}
\begin{tikzpicture}
\begin{axis}[bkplot, width=86mm, height=46mm, xmin=0, xmax=65, ymin=0, ymax=90, xtick={0,20,40,60}, ytick={20,50,80}, xlabel={$t$（分）}, ylabel={$T$}]
  \addplot[bkthin, dashed, domain=0:62] {20};
  \addplot[bkcurve, domain=0:62, samples=50] {20+60*exp(-0.05*x)};
  \draw[bkOrange, line width=1pt, <->] (axis cs:20,20.8) -- (axis cs:20,41.3);
  \node[font=\sffamily\scriptsize, anchor=east, text=bkOrange!80!black] at (axis cs:18.3,30) {差 $u$};
  \node[font=\sffamily\scriptsize, anchor=north east, text=bkGray] at (axis cs:62,16) {部屋 20\,℃};
\end{axis}
\end{tikzpicture}
\bkfigcaption{お湯の温度 $T$（青）と、部屋の温度との差 $u$（オレンジの矢印）。$T$ と $u$ はいつも20ちがうだけ。}
\end{bkfigure}
$T$ と $u$ は、いつも20ちがうだけなので、$u$ のグラフは $T$ のグラフを20だけ下にずらしたものです。ずらしても傾きは変わりません。式では、20は定数で微分すると0なので、$u'=T'-0=T'$ です。
$T'=-k(T-20)$ の左の $T'$ を $u'$ に、右の $T-20$ を $u$ にかえると、
\[ u'=-ku \]
これは表の「微分すると自分の $a$ 倍」（$a=-k$）なので、答えは $u=Ce^{-kt}$。$u=T-20$ をもどすと、
\[ T-20=Ce^{-kt},\qquad\text{つまり}\qquad T=20+Ce^{-kt} \]
「差」が $e^{-kt}$ でへり、温度は20\,℃に近づきます（20\,℃より下にはなりません）。5-16の置換と同じ「おきかえ」の考えです。

\bksubsection{5}{ふるまいをくらべる：どれがいちばん速くふえる？}

\label{bk:5-18:kurabe}$t$ が大きくなったときのふえ方をくらべます。
\begin{bkcenter}
\begin{tblr}{colspec={Q[c,m]*{7}{Q[r,m]}}, row{1}={bg=bkTheme!12, font=\sffamily\small}, hline{1,2,Z}={0.6pt,bkGray}, colsep=5pt, cells={font=\small}}
$t$ & 1 & 2 & 3 & 4 & 5 & 10 & 20 \\
$\ln t$ & 0 & 0.69 & 1.10 & 1.39 & 1.61 & 2.30 & 3.00 \\
$t^3$ & 1 & 8 & {\bfseries 27} & {\bfseries 64} & 125 & 1000 & 8000 \\
$e^t$ & 2.72 & 7.39 & 20.1 & 54.6 & {\bfseries 148} & {\bfseries 22026} & {\bfseries 485165195} \\
\end{tblr}
\end{bkcenter}
表をよく見ましょう。$t=3$、4では $t^3$ のほうが大きいのですが、$t=5$ で $e^t$ が追いこし、そのあとは、どんどん引きはなします。$t=20$ では、$e^t$ は $t^3$ の6万倍以上です。

なぜ、最後は $e^t$ が勝つのでしょう。5-17のテイラー展開を思い出します。
\[ e^t=1+t+\frac{t^2}{2!}+\frac{t^3}{3!}+\frac{t^4}{4!}+\frac{t^5}{5!}+\cdots \]
$e^t$ の中には、$t^3$ より次数の高い $\dfrac{t^4}{24}$ がふくまれています。$\dfrac{t^4}{24}$ は $t^3$ の $\dfrac{t}{24}$ 倍なので、$t$ が24より大きくなれば、この1つの項だけで $t^3$ をこえます。そのうえ $t^5$、$t^6$、… の項もあるので、いずれ必ず $e^t$ が勝つのです。逆に $\ln t$ はとてもゆっくりふえます。$t$ が10倍になっても、$\ln t$ は2.3ふえるだけです（$\ln10=2.30$）。$t$ が大きいとき、
\[ \ln t\ \ll\ t^n\ \ll\ e^t\qquad\text{（$\ll$ は「ずっと小さい」）} \]
同じ理由で、$t^3e^{-t}=\dfrac{t^3}{e^t}$ は、$t$ が大きくなると0に近づきます（$t=10$ で0.045、$t=20$ で0.00002）。「$e^{-t}$ がかかると、$t^n$ がかかっていても、やがて0に近づく」ことは、9-5のラプラス変換でも使います。

\begin{bkfigure}
\begin{tikzpicture}[x=1mm, y=1mm, nd/.style={draw=bkTheme, line width=0.8pt, rounded corners=3pt, fill=bkTheme!6, font=\small, align=center, inner sep=3pt, minimum height=13mm, text width=26mm}]
  \node[nd] (a) at (0,14) {平均の速さ\\（5-1）};
  \node[nd] (b) at (36,14) {微分\\（5-3〜5-8）};
  \node[nd] (c) at (72,14) {積分（微分の逆）\\（5-9〜5-13）};
  \node[nd, fill=bkOrange!12, draw=bkOrange] (d) at (108,14) {$e$・$\sin$・$\cos$\\（5-14〜5-17）};
  \node[nd, fill=bkPurple!10, draw=bkPurple, text width=108mm] (e) at (54,-8) {第6巻：微分方程式\enspace「変わり方のきまり」から未来を求める};
  \draw[bk arrow, bkTheme, line width=1pt, shorten <=1.5mm, shorten >=1.5mm] (a) -- (b);
  \draw[bk arrow, bkTheme, line width=1pt, shorten <=1.5mm, shorten >=1.5mm] (b) -- (c);
  \draw[bk arrow, bkTheme, line width=1pt, shorten <=1.5mm, shorten >=1.5mm] (c) -- (d);
  \draw[bk arrow, bkPurple, line width=1.2pt, shorten <=1.5mm, shorten >=1.5mm] (d.south) -- (e.north -| d.south);
  \draw[bk arrow, bkPurple, line width=1.2pt, shorten <=1.5mm, shorten >=1.5mm] (b.south) -- (e.north -| b.south);
\end{tikzpicture}%
\bkfigcaption{第5巻の地図。微分と積分の道具と、答えの候補の関数をもって、第6巻へ。}
\end{bkfigure}

\Needspace{12\baselineskip}
\bksubsection{6}{第5巻の計算のきまり 一覧}

\label{bk:5-18:kimari}第6巻からは、ここにあるきまりを、ことわりなしに使います。わからなくなったら、右の章にもどりましょう。
\begin{bkcenter}
\small
\begin{tblr}{width=\linewidth, colspec={Q[l,m]X[l,m]Q[c,m]}, row{1}={bg=bkTheme!12, font=\sffamily\small}, hline{1,2,Z}={0.6pt,bkGray}, hline{3-12}={0.3pt,bkGray}, rowsep=2.5pt, colsep=5pt, column{1}={font=\sffamily\small}}
きまり & 式 & 章 \\
微分の定義 & $f'(x)=\displaystyle\lim_{h\to0}\frac{f(x+h)-f(x)}{h}$（接線の傾き・瞬間の変わる速さ） & 5-3 \\
たし算・定数倍 & $(f+g)'=f'+g'$、$(cf)'=cf'$ & 5-4 \\
積・合成関数 & $(fg)'=f'g+fg'$、$\dfrac{\dd y}{\dd x}=\dfrac{\dd y}{\dd u}\cdot\dfrac{\dd u}{\dd x}$ & 5-6 \\
増減と最大・最小 & $f'>0$ でふえる、$f'<0$ でへる。$f'=0$ の所が山や谷の候補 & 5-7 \\
1次近似 & $f(a+h)\approx f(a)+f'(a)\,h$ & 5-8 \\
定積分 & 細かく分けた長方形の面積の和 $\displaystyle\int_a^bf\,\dd x=\lim_{\Delta x\to0}\sum f(x_k)\,\Delta x$ & 5-10 \\
微積分の基本定理 & $F'=f$ なら $\displaystyle\int_a^bf\,\dd x=F(b)-F(a)$ & 5-11 \\
初期条件 & 積分の $+C$ は、はじめの状態で決まる & 5-13 \\
置換・部分積分 & $\dd u=u'\,\dd x$ で書きかえ。$\displaystyle\int fg'\,\dd x=fg-\int f'g\,\dd x$ & 5-16 \\
テイラー展開 & $f(x)=f(0)+f'(0)x+\dfrac{f''(0)}{2!}x^2+\cdots$ & 5-17 \\
\end{tblr}
\end{bkcenter}

\begin{motto}[4]
\textbf{$f'=af$ の答えは $Ce^{at}$ だけ}\par
ほかに答えはないのでしょうか。$f'=af$ をみたす $f$ について、$g=f\,e^{-at}$ を積の微分（5-6）で微分すると、
\[ g'=f'e^{-at}+f\cdot(-a)e^{-at}=af\,e^{-at}-af\,e^{-at}=0 \]
（$f'=af$ を入れました。）微分して0になるのは定数なので、$g=C$。$g=fe^{-at}$ の両辺に $e^{at}$ をかけると、$e^{-at}\times e^{at}=e^0=1$ なので $f=Ce^{at}$。答えはこれだけだと分かります。
\end{motto}

\begin{senmonsho}[大学の教科書ふうの書き方]
\begin{bkhonbun}
微分方程式 $\dfrac{\dd y}{\dd t}=ay$ の一般解は $y=Ce^{at}$（$C$ は任意定数）であり、初期条件 $y(0)=y_0$ のもとで特殊解 $y=y_0e^{at}$ が一意に定まる。\bkidx{いっぱんかい}{一般解}\bkidx{とくしゅかい}{特殊解}
\end{bkhonbun}
\bkyaku{「一般解」は $C$ がついたままの答え、「任意定数」はどんな数でもよい定数 $C$、「特殊解」は $C$ を決めた答え。「一意に定まる」は「ただ1つに決まる」。}
\end{senmonsho}

\begin{taisetsu}
\def\tsl#1{\makebox[5.2\zw][l]{\sffamily\bfseries\small\color{bkTheme}#1}}
\tsl{答えの候補}$f'=af$ → $Ce^{at}$\qquad $f''=-\omega^2f$ → $A\sin\omega t+B\cos\omega t$\\[3pt]
\tsl{}$f'=v$ → $vt+C$\qquad $f''=a^2f$ → $Ae^{at}+Be^{-at}$\\[3pt]
\tsl{数の決め方}$C$、$A$、$B$ は、はじめの状態（$t=0$ の値）で決める。\\[3pt]
\tsl{ふえ方}$\ln t\ \ll\ t^n\ \ll\ e^t$
\end{taisetsu}

\begin{kotoba}
\small
\bkword{線形}{答えどうしをたしても、答えを何倍しても答えになる式の性質。}
\bkword{初期条件}{はじめ（$t=0$）の位置や速度。式の中の数を決める。}
\end{kotoba}

%% ---------------- 例題 ----------------
\bksection{reidai}{例題でたしかめよう}

\begin{reidai}[自分に比例してへる]
$f'=-3f$、$f(0)=5$ をみたす関数を求めましょう。また、$t=1$ での値を求めましょう（$e^{-3}=0.0498$）。\par
\kangaekata 「微分すると自分の $-3$ 倍」なので、②の表から答えの形は $Ce^{-3t}$。$C$ は $t=0$ の値で決める（③）。\par
\shiki 答えの形は $f=Ce^{-3t}$。$t=0$ で $f(0)=Ce^0=C=5$\par
\hspace*{2.6\zw}確かめ：$f'=5\times(-3)e^{-3t}=-3\times5e^{-3t}=-3f$\par
\hspace*{2.6\zw}$t=1$ で $5e^{-3}=5\times0.0498=0.249$\par
\kotae $f=5e^{-3t}$、$t=1$ で約0.25
\end{reidai}

\begin{jibun}
$f'=0.5f$、$f(0)=2$ をみたす関数を求めましょう。$t=2$ での値も求めましょう（$e=2.718$）。
\bkjans{$f=2e^{0.5t}$。$t=2$ で $2e^{1}=2e\approx5.44$}
\end{jibun}

\begin{reidai}[ゆれる動きを決める]
$\ddot{x}=-9x$、$x(0)=0.2$、$\dot{x}(0)=0$ をみたす $x$ を求めましょう。また、周期を求めましょう。\par
\kangaekata 「2回微分で自分の $-9$ 倍」なので $\omega^2=9$、$\omega=3$（②）。$A$、$B$ ははじめの位置と速度で決める（③）。周期は $\frac{2\pi}{\omega}$（5-15）。\par
\shiki 答えの形は $x=A\sin3t+B\cos3t$。\par
\hspace*{2.6\zw}$x(0)=A\times0+B\times1=B=0.2$\par
\hspace*{2.6\zw}$\dot{x}=3A\cos3t-3B\sin3t$ より $\dot{x}(0)=3A=0$、$A=0$\par
\hspace*{2.6\zw}周期 $\dfrac{2\pi}{3}=\dfrac{6.28}{3}\approx2.09$\par
\kotae $x=0.2\cos3t$、周期約2.1
\end{reidai}

\begin{jibun}
$\ddot{x}=-16x$、$x(0)=0$、$\dot{x}(0)=2$ をみたす $x$ を求めましょう。周期も求めましょう。
\bkjans{$x=A\sin4t+B\cos4t$。$B=0$、$4A=2$ より $x=\dfrac{1}{2}\sin4t$。周期 $\dfrac{2\pi}{4}\approx1.57$}
\end{jibun}

\begin{reidai}[位置も速度も0でないとき]
$\ddot{x}=-4x$、$x(0)=0.1$、$\dot{x}(0)=0.6$ をみたす $x$ を求めましょう。\par
\kangaekata $\omega^2=4$ より $\omega=2$。はじめの位置で $B$ を、はじめの速度で $A$ を決める（③の例2と同じ手順）。\par
\shiki 答えの形は $x=A\sin2t+B\cos2t$\par
\hspace*{2.6\zw}位置：$x(0)=A\times0+B\times1=B$。$B=0.1$\par
\hspace*{2.6\zw}速度：$\dot{x}=2A\cos2t-2B\sin2t$ より\par
\hspace*{2.6\zw}$\dot{x}(0)=2A\times1-2B\times0=2A$\par
\hspace*{4.6\zw}$2A=0.6$ より $A=0.3$\par
\kotae $x=0.3\sin2t+0.1\cos2t$
\end{reidai}
$\sin$ の項は速度で、$\cos$ の項は位置で決まります。$t=0$ で $\sin0=0$ なので、$\sin$ の項は位置に効かず、微分すると $\cos$ になるので速度に効くからです。

\begin{jibun}
$\ddot{x}=-9x$、$x(0)=0.2$、$\dot{x}(0)=0.3$ をみたす $x$ を求めましょう。
\bkjans{$x=A\sin3t+B\cos3t$。$B=0.2$、$3A=0.3$ より $A=0.1$。$x=0.1\sin3t+0.2\cos3t$}
\end{jibun}

\begin{reidai}[式をみたすか確かめる]
$y=3e^{2t}+e^{-2t}$ が $y''=4y$ をみたすことを確かめましょう。\par
\kangaekata ①の表で2回微分し、$y''$ が $y$ の4倍になるかを見る。\par
\shiki $y'=3\times2e^{2t}+(-2)e^{-2t}=6e^{2t}-2e^{-2t}$\par
\hspace*{2.6\zw}$y''=6\times2e^{2t}-2\times(-2)e^{-2t}=12e^{2t}+4e^{-2t}$\par
\hspace*{2.6\zw}$12e^{2t}+4e^{-2t}=4(3e^{2t}+e^{-2t})=4y$\par
\kotae $y''=4y$ をみたす
\end{reidai}

\begin{jibun}
$y=2\sin3t-\cos3t$ が $y''=-9y$ をみたすことを確かめましょう。
\bkjans{$y'=6\cos3t+3\sin3t$、$y''=-18\sin3t+9\cos3t=-9(2\sin3t-\cos3t)=-9y$}
\end{jibun}

\begin{reidai}[冷めるお湯]
部屋の温度は20\,℃です。お湯の温度 $T$ が、$T'=-0.05(T-20)$ にしたがって冷めます（$t$ は分）。はじめ80\,℃のとき、$T$ の式と20分後の温度を求めましょう。また、50\,℃になるのは何分後ですか（$e^{-1}=0.368$、$\ln2=0.693$）。\par
\kangaekata ④のように、部屋との温度差 $T-20$ が $e^{-0.05t}$ でへる。時刻は、両辺の自然対数をとって求める（5-14）。\par
\shiki $T=20+Ce^{-0.05t}$。$t=0$ で $20+C\times1=80$ より $C=60$\par
\hspace*{2.6\zw}20分後：$T=20+60e^{-0.05\times20}=20+60e^{-1}$\par
\hspace*{4.6\zw}$=20+60\times0.368=20+22.1$\par
\hspace*{2.6\zw}50\,℃：$20+60e^{-0.05t}=50$ より $60e^{-0.05t}=30$、$e^{-0.05t}=\dfrac{1}{2}$\par
\hspace*{2.6\zw}両辺の $\ln$ をとる：$-0.05t=\ln\dfrac{1}{2}=-\ln2$\par
\hspace*{4.6\zw}$t=\dfrac{\ln2}{0.05}=\dfrac{0.693}{0.05}=13.9$\par
\kotae $T=20+60e^{-0.05t}$、約42\,℃、約14分後
\end{reidai}
$\ln\dfrac{1}{2}=-\ln2$ は、$e^{-0.693}=\dfrac{1}{e^{0.693}}=\dfrac{1}{2}$ だからです（3-12の $\log_ab^r=r\log_ab$ で $r=-1$）。50\,℃は、はじめの差60\,℃が半分の30\,℃になった時刻です。

\begin{jibun}
部屋の温度25\,℃で、$T'=-0.1(T-25)$（$t$ は分）、はじめ90\,℃のお湯の、10分後の温度を求めましょう（$e^{-1}=0.368$）。
\bkjans{$T=25+65e^{-0.1t}$。$t=10$ で $25+65\times0.368\approx48.9\,℃$}
\end{jibun}

\begin{reidai}[表を使って微分・積分する]
$f(t)=3t^2+2e^{-t}+\sin2t$ を微分しましょう。また、$\displaystyle\int_0^{\pi}f(t)\,\dd t$ を求めましょう（$\pi^3=31.0$、$e^{-\pi}=0.043$）。\par
\kangaekata 項ごとに①の表を使う（定数倍は前に出す、5-4）。定積分は上端と下端の値の差（5-11）。\par
\shiki $f'=3\times2t+2\times(-e^{-t})+2\cos2t=6t-2e^{-t}+2\cos2t$\par
\hspace*{2.6\zw}原始関数：$3\times\dfrac{t^3}{3}+2\times(-e^{-t})-\dfrac{1}{2}\cos2t=t^3-2e^{-t}-\dfrac{1}{2}\cos2t$\par
\hspace*{2.6\zw}$\begin{aligned}[t]\int_0^{\pi}f(t)\,\dd t&=\Bigl(\pi^3-2e^{-\pi}-\frac{1}{2}\cos2\pi\Bigr)-\Bigl(0-2e^0-\frac{1}{2}\cos0\Bigr)\\
&=\Bigl(\pi^3-2e^{-\pi}-\frac{1}{2}\Bigr)-\Bigl(-2-\frac{1}{2}\Bigr)=\pi^3+2-2e^{-\pi}\end{aligned}$\par
\kotae $f'=6t-2e^{-t}+2\cos2t$、積分は約 $31.0+2-0.09=32.9$
\end{reidai}
$\sin2t$ の部分は、
\[ \Bigl(-\frac{1}{2}\cos2\pi\Bigr)-\Bigl(-\frac{1}{2}\cos0\Bigr)=-\frac{1}{2}+\frac{1}{2}=0 \]
で、答えに残りませんでした。0から $\pi$ までで $\sin2t$ はちょうど1周期なので、山と谷の面積が打ち消しあうからです（5-15）。このように、項ごとに分けて計算すると、どの項がどれだけ答えにきいているかも分かります。

\begin{jibun}
$g(t)=4t^3-e^{2t}+3\cos t$ を微分しましょう。また、$\displaystyle\int_0^1g(t)\,\dd t$ を求めましょう（$e^2=7.389$、$\sin1=0.841$）。
\bkjans{$g'=12t^2-2e^{2t}-3\sin t$。$\Bigl[\,t^4-\dfrac{1}{2}e^{2t}+3\sin t\,\Bigr]_0^1=1-\dfrac{e^2}{2}+3\sin1+\dfrac{1}{2}\approx0.33$}
\end{jibun}

\begin{reidai}[たおれる動き]
少しかたむいた棒のかたむき $\theta$（rad）が、$\ddot{\theta}=4\theta$ にしたがって大きくなります。$\theta(0)=0.01$、$\dot{\theta}(0)=0$ のとき、$\theta$ の式と、2秒後のかたむきを求めましょう（$e^4=54.6$、$e^{-4}=0.018$）。\par
\kangaekata 「2回微分で自分の4倍」なので、②の表で $a^2=4$、$a=2$。答えの形は $Ae^{2t}+Be^{-2t}$。$A$、$B$ ははじめの状態で決める（③）。\par
\shiki $\theta=Ae^{2t}+Be^{-2t}$。$t=0$ で $e^0=1$ なので $\theta(0)=A+B=0.01$\par
\hspace*{2.6\zw}$\dot{\theta}=2Ae^{2t}-2Be^{-2t}$ より $\dot{\theta}(0)=2A-2B=0$、つまり $A=B$\par
\hspace*{2.6\zw}$A+B=0.01$、$A=B$ より $A=B=0.005$\par
\hspace*{2.6\zw}2秒後：$0.005\times(e^4+e^{-4})=0.005\times(54.6+0.018)\approx0.273$\par
\kotae $\theta=0.005(e^{2t}+e^{-2t})$、約0.27\,rad（はじめの27倍）
\end{reidai}
$e^{-2t}$ の部分はすぐに小さくなり、$e^{2t}$ の部分がどんどん大きくなります。はじめのかたむきがほんの少しでも、たおれる動きは止まらないことが分かります。

\begin{jibun}
$\ddot{x}=9x$、$x(0)=0$、$\dot{x}(0)=0.6$ をみたす $x$ を求め、$t=1$ での値を求めましょう（$e^3=20.09$、$e^{-3}=0.0498$）。
\bkjans{$x=Ae^{3t}+Be^{-3t}$。$A+B=0$、$3A-3B=0.6$ より $A=0.1$、$B=-0.1$。$x=0.1(e^{3t}-e^{-3t})$、$t=1$ で約2.00}
\end{jibun}

%% ---------------- まちがえやすいところ・物理工学 ----------------
\bksection{tsumazuki}{まちがえやすいところをチェック}

\begin{tsumazuki}{$\omega$ と $\omega^2$ を取りちがえる}
\bkng{$\ddot{x}=-9x$ の答えは $\sin9t$}
\bkok{$\ddot{x}=-9x$ の答えは $\sin3t$}
2回微分で前に出るのは $\omega^2$ なので、$\omega^2=9$ から $\omega=3$ です。$\sin9t$ を2回微分すると $-81\sin9t$ で、$-9$ 倍になりません。
\end{tsumazuki}

\begin{tsumazuki}{はじめの状態を、とちゅうで使う}
$C$ を決めるのは $t=0$ の値です。$T=20+Ce^{-0.05t}$ で $T(0)=80$ なら、$e^0=1$ なので $20+C=80$。$t$ に0を入れてから $C$ を求めます。
\end{tsumazuki}

\begin{tsumazuki}{2回微分の式で、数を1つしか決めない}
$\ddot{x}=-4x$ の答えの形は $A\sin2t+B\cos2t$ で、数が2つあります。はじめの位置だけでは $B$ しか決まりません。はじめの速度も使って $A$ を決めます。
\end{tsumazuki}

\begin{butsuri}
\begin{itemize}[leftmargin=1.2\zw, itemsep=1pt, topsep=0pt, label={\color{bkPurple}●}]
\item \textbf{ふえる・へる（$e^{at}$）}\par\noindent 冷めるお湯、コンデンサーの放電、薬の濃度、人口。答えはいつも $Ce^{at}$ の形（6-4）。
\item \textbf{ゆれる（$\sin$・$\cos$）}\par\noindent ばね、振り子、コイルとコンデンサーをつないだ回路（LC回路）の電気の振動、建物のゆれ（6-6）。
\item \textbf{ゆれながら止まる（$e^{at}$ と $\sin$ のかけ算）}\par\noindent 計器の針、車のサスペンション（6-7・6-12）。2つの候補をかけ算した形が答えになります。
\item \textbf{制御工学（第9巻）}\par\noindent ラプラス変換の表（9-6）にならぶのも、この章の表と同じ関数です。
\end{itemize}
\end{butsuri}

%% ---------------- 章末問題 ----------------
\bkendproblems
\bkproblemgroup{基本}{3}

\bkproblem 次のなぞなぞに合う関数を、1つずつ答えましょう。\newline
（1）微分すると、自分の $-4$ 倍\quad（2）2回微分すると、自分の $-25$ 倍\newline
（3）微分すると、いつも $-3$
\bkpans{（1）\nobreak $Ce^{-4t}$\quad（2）\nobreak $A\sin5t+B\cos5t$\quad（3）\nobreak $-3t+C$}

\bkproblem $f'=-0.2f$、$f(0)=50$ をみたす関数を求め、$t=5$ での値を求めましょう（$e^{-1}=0.368$）。
\bkpans{$f=50e^{-0.2t}$。$t=5$ で $50e^{-1}\approx18.4$}

\bkproblem $\ddot{x}=-25x$、$x(0)=0.05$、$\dot{x}(0)=0$ をみたす $x$ を求めましょう。周期も求めましょう。
\bkpans{$x=0.05\cos5t$。周期 $\dfrac{2\pi}{5}\approx1.26$}

\bkproblem $y=e^{3t}-2e^{-3t}$ が $y''=9y$ をみたすことを確かめましょう。
\bkpans{$y''=9e^{3t}-18e^{-3t}=9(e^{3t}-2e^{-3t})=9y$}

\bkproblem $h(t)=2t^3-3e^{2t}+4\sin t$ を微分しましょう。また、$\displaystyle\int_0^{\pi}4\sin t\,\dd t$ を求めましょう。
\bkpans{$h'=6t^2-6e^{2t}+4\cos t$。$\Bigl[\,-4\cos t\,\Bigr]_0^{\pi}=4+4=8$}

\bkproblemgroup{チャレンジ}{3}

\bkproblem 部屋の温度22\,℃で、$T'=-0.04(T-22)$（$t$ は分）にしたがって冷めるコーヒーがあります。はじめ82\,℃です。$T$ の式を求め、25分後の温度と、温度が50\,℃になる時刻を求めましょう（$e^{-1}=0.368$、$\ln\frac{15}{7}=0.762$）。
\bkpans{$T=22+60e^{-0.04t}$。25分後 $22+60\times0.368\approx44\,℃$。$60e^{-0.04t}=28$ より \mbox{$t=\dfrac{\ln\frac{60}{28}}{0.04}=\dfrac{0.762}{0.04}\approx19$\,分後}}

\bkproblem $\ddot{x}=-4x$、$x(0)=0.3$、$\dot{x}(0)=0.8$ をみたす $x$ を求めましょう。
\bkpans{$x=A\sin2t+B\cos2t$。$B=0.3$、$2A=0.8$ より $A=0.4$。$x=0.4\sin2t+0.3\cos2t$}

\bkproblem 部屋の温度20\,℃の部屋に、5\,℃の冷たい飲み物をおきました。温度 $T$ は $T'=-0.1(T-20)$（$t$ は分）にしたがってあたたまります。$T$ の式と、10分後の温度を求めましょう（$e^{-1}=0.368$）。
\bkpans{$T=20+Ce^{-0.1t}$。$t=0$ で $20+C=5$ より $C=-15$。$T=20-15e^{-0.1t}$。10分後 $20-15\times0.368\approx14.5\,℃$}

\bkproblem $\ddot{x}=x$、$x(0)=2$、$\dot{x}(0)=0$ をみたす $x$ を求めましょう。
\bkpans{$x=Ae^{t}+Be^{-t}$。$A+B=2$、$A-B=0$ より $A=B=1$。$x=e^t+e^{-t}$}

\begin{furikaeri}
\begin{checklist}
\item ふるまい表を使って、微分・積分できる。\hfill{\footnotesize\sffamily\color{bkTheme}→\pageref{bk:5-18:hyo}ページ}
\item 「自分にもどる」関数を答えられる。\hfill{\footnotesize\sffamily\color{bkTheme}→\pageref{bk:5-18:modoru}ページ}
\item はじめの状態から、式の中の数を決められる。\hfill{\footnotesize\sffamily\color{bkTheme}→\pageref{bk:5-18:shoki}ページ}
\end{checklist}
\end{furikaeri}

\bkfukidashi{delta}{第5巻、おつかれさま！　答えの候補のカードはそろった。第6巻では、これを使って、ものの動きを式で求めよう。}
